% -----------------------------------------------------------------------------
% File: sections/04_llm_code_review.tex
% Description: Section 4 - LLMs as Code Review Agents
% -----------------------------------------------------------------------------

\section{Paper 3: LLMs as Code Review Agents}
\label{sec:llm_review}

\paperCard{LLMs as Code Review Agents: A Rapid Review and Experimental Evaluation with Human Expert Judges}{M. Kawalerowicz, M. Pietranik, K. St\k{e}pniak}{ICCCI 2025}

\subsection{Background and Research Motivation}
Quality assurance through code review is a vital phase of the software development lifecycle (SDLC), ensuring that code meets functional requirements, adheres to design principles, and is free from security vulnerabilities. However, manual code review by human developers is resource-intensive, time-consuming, and often becomes a bottleneck that slows down development velocity.

Traditional automated tools such as static analyzers and linters are limited to syntactic rules and pattern matching, failing to assess higher-level semantic properties like code readability, adherence to SOLID design principles (SRP, OCP, LSP, ISP, DIP), or the identification of complex security risks and architectural anti-patterns.

Recently, Large Language Models (LLMs) have been proposed as automated code review agents under the \highlight{LLM-as-a-Judge} paradigm. While LLMs are significantly cheaper and faster than human reviewers, assessing their alignment with human expert consensus and understanding their biases is crucial before deploying them in production code review pipelines. Key questions include: How well do LLMs correlate with senior developer evaluations? Which metrics do they excel at? And what is the cost-benefit trade-off between accuracy and API expenses?

\subsection{Paper Objectives}
The paper conducts a comparative evaluation of LLM-generated code quality assessments against human expert judgments. The specific objectives include:
\begin{itemize}
    \item Designing a parallel evaluation framework where both LLM agents and human experts independently score the same code corpus.
    \item Evaluating 6 LLMs (Claude 3.5 Sonnet, Claude 3.5 Haiku, GPT-4o-mini, Mistral Large, Ministral-8b, Falcon3-7B) against 3 human software developers with 5--15 years of industry experience.
    \item Measuring alignment using statistically rigorous non-parametric methods (Shapiro-Wilk test, Spearman rank correlation).
    \item Providing a cost-benefit analysis to determine which LLM offers the best accuracy-to-cost ratio for production deployment.
\end{itemize}

\subsection{Proposed Method}

\subsubsection{Parallel Evaluation Framework}
The study establishes a controlled experimental framework where both human experts and LLM agents independently evaluate the same set of code files. A custom evaluation platform is developed for human reviewers, while LLMs receive structured prompts with identical evaluation criteria.

\begin{figure}[h]
    \centering
    \resizebox{0.95\textwidth}{!}{%
    \begin{tikzpicture}[
        node distance=0.5cm and 1.0cm,
        font=\sffamily\small,
        corpus/.style={rectangle, draw=codeGreen, fill=green!5, thick, minimum width=2.5cm, minimum height=1.2cm, align=center, rounded corners},
        actor/.style={rectangle, draw=black, thick, fill=white, minimum width=3cm, minimum height=1.0cm, align=center, drop shadow},
        process/.style={rectangle, draw=gray, dashed, fill=gray!5, minimum width=2.5cm, minimum height=0.8cm, align=center, font=\footnotesize},
        metric/.style={circle, draw=black, fill=yellow!10, thick, inner sep=2pt, minimum size=0.8cm},
        arrow/.style={->, >=stealth, ultra thick, gray!80, rounded corners=5pt}
    ]
        \node[corpus] (input) {\textbf{Code Corpus ($F$)}\\75 files\\(Java, Py, C\#, JS, TS)};
        \node[actor, above right=0.5cm and 2.0cm of input, draw=techBlue, fill=blue!5] (human) {\textbf{Human Experts}\\($n=3$ Seniors)};
        \node[actor, below right=0.5cm and 2.0cm of input, draw=orange, fill=orange!5] (llm) {\textbf{LLM Agents}\\(GPT, Claude, Mistral)};
        \node[process, right=0.8cm of human] (platform) {Custom Platform};
        \node[process, right=0.8cm of llm] (prompt) {Structured Prompts};
        \coordinate (midpoint) at ($(platform)!0.5!(prompt)$);
        \node[metric, right=2.0cm of midpoint, align=center] (metrics) {$\vec{m}$};
        \node[below=0.1cm of metrics, font=\scriptsize, align=center] {12 Metrics\\($m_1 \dots m_{12}$)};
        \node[corpus, right=1.5cm of metrics, fill=gray!10, draw=black] (output) {\textbf{Score Matrices}\\$S_{Human} \text{ vs } S_{LLM}$};
        \draw[arrow] (input.east) -- ++(0.5,0) |- (human.west);
        \draw[arrow] (input.east) -- ++(0.5,0) |- (llm.west);
        \draw[arrow] (human) -- (platform);
        \draw[arrow] (llm) -- (prompt);
        \draw[arrow] (platform.east) -- ++(0.5,0) |- (metrics.west);
        \draw[arrow] (prompt.east) -- ++(0.5,0) |- (metrics.west);
        \draw[arrow] (metrics) -- (output);
    \end{tikzpicture}%
    }
    \caption{Methodology flow of the parallel code quality evaluation framework.}
    \label{fig:llm_flow}
\end{figure}

\subsubsection{Code Corpus and Evaluators}
\begin{itemize}
    \item \textbf{Dataset:} 75 source code files in Java, C\#, Python, JavaScript, and TypeScript, selected to represent varying levels of code quality from poor to excellent.
    \item \textbf{Human evaluators:} 3 senior software developers with 5--15 years of industry experience, each independently scoring all 75 files.
    \item \textbf{LLM evaluators:} 6 LLMs of varying sizes and capabilities: Claude 3.5 Sonnet, Claude 3.5 Haiku, GPT-4o-mini, Mistral Large, Ministral-8b, and Falcon3-7B.
\end{itemize}

\subsubsection{Code Quality Metrics}
Each code file is scored on a scale of 1--10 across 12 quality metrics:
\begin{itemize}
    \item \textbf{Comments} ($m_1$), \textbf{DIP} ($m_2$), \textbf{Formatting} ($m_3$), \textbf{ISP} ($m_4$), \textbf{KISS} ($m_5$), \textbf{LSP} ($m_6$)
    \item \textbf{Naming} ($m_7$), \textbf{OCP} ($m_8$), \textbf{Performance} ($m_9$), \textbf{DRY/Repetition} ($m_{10}$), \textbf{Safety} ($m_{11}$), \textbf{SRP} ($m_{12}$)
\end{itemize}

\subsubsection{Statistical Analysis Pipeline}
The study establishes a rigorous statistical pipeline to validate its findings:
\begin{enumerate}
    \item \textbf{Normality Check:} The Shapiro-Wilk test is performed on all score distributions. The resulting $p$-values are less than $0.05$ ($p < \alpha$), indicating that scores do not follow a normal distribution.
    \item \textbf{Correlation Metric:} Because the data is non-normal, the non-parametric \highlight{Spearman Rank Correlation} ($\rho$) is used:
    \begin{equation}
        \rho = 1 - \frac{6 \sum d_i^2}{n(n^2 - 1)}
    \end{equation}
    where $d_i$ is the difference between ranks.
    \item \textbf{Ground Truth Modeling:} Human developer consensus is established as the ground truth baseline:
    \begin{equation}
        GT(f_i) = \text{Median}(S_{Human1}(f_i), S_{Human2}(f_i), S_{Human3}(f_i))
    \end{equation}
\end{enumerate}

\subsection{Experiments and Results}

\subsubsection{Spearman Correlation with Human Consensus}
Table~\ref{tab:llm_metrics} presents the Spearman correlation of each LLM with the human consensus baseline across all 12 quality metrics.

\begin{table}[h]
    \centering
    \caption{Spearman correlation ($\rho$) of LLM agents with human consensus.}
    \label{tab:llm_metrics}
    \resizebox{\linewidth}{!}{%
    \begin{tabular}{lcccccc}
        \toprule
        \textbf{Metric} & \textbf{Sonnet} & \textbf{GPT-4o-m} & \textbf{Mistral L} & \textbf{Haiku} & \textbf{Min-8b} \\
        \midrule
        $m_1$ (Comm.) & \textbf{0.71} & 0.64 & 0.57 & 0.62 & 0.43 \\
        $m_2$ (DIP)   & \textbf{0.83} & 0.58 & 0.54 & 0.50 & 0.30 \\
        $m_3$ (Format)& \textbf{0.71} & 0.66 & 0.60 & 0.65 & 0.48 \\
        $m_4$ (ISP)   & \textbf{0.90} & 0.81 & 0.76 & 0.44 & 0.28 \\
        $m_5$ (KISS)  & \textbf{0.85} & 0.76 & 0.79 & 0.66 & 0.59 \\
        $m_6$ (LSP)   & \textbf{0.88} & 0.74 & 0.69 & 0.28 & 0.17 \\
        $m_7$ (Name)  & 0.80 & 0.72 & \textbf{0.81} & 0.70 & 0.59 \\
        $m_8$ (OCP)   & \textbf{0.79} & 0.65 & 0.73 & 0.77 & 0.57 \\
        $m_9$ (Perf)  & 0.73 & 0.66 & \textbf{0.73} & 0.53 & 0.58 \\
        $m_{10}$ (DRY)& \textbf{0.87} & 0.79 & 0.87 & 0.78 & 0.59 \\
        $m_{11}$ (Safe)& 0.76 & 0.75 & \textbf{0.77} & 0.70 & 0.69 \\
        $m_{12}$ (SRP) & \textbf{0.86} & 0.78 & 0.85 & 0.71 & 0.63 \\
        \bottomrule
    \end{tabular}
    }
\end{table}

Detailed analysis of the results:
\begin{itemize}
    \item \textbf{Claude 3.5 Sonnet} achieves the highest alignment with human consensus in \highlight{9 out of 12 metrics}, demonstrating particularly strong performance on complex SOLID design principles: ISP ($\rho = 0.90$), LSP ($\rho = 0.88$), DRY ($\rho = 0.87$), and SRP ($\rho = 0.86$).
    \item \textbf{Mistral Large} achieves the best scores in 3 metrics: Naming ($\rho = 0.81$), Performance ($\rho = 0.73$), and Safety ($\rho = 0.77$), suggesting strength in surface-level code evaluation.
    \item \textbf{Smaller models} (Ministral-8b, Falcon3-7B) show significantly lower correlations across all metrics, particularly in complex design principles like LSP ($\rho = 0.17$) and ISP ($\rho = 0.28$), indicating that model size plays a crucial role in understanding abstract software engineering concepts.
    \item \textbf{GPT-4o-mini} provides consistently moderate correlations ($\rho \approx 0.65$--$0.81$) across all metrics, positioning it as a balanced mid-tier option.
\end{itemize}

\subsubsection{Cost-Benefit Analysis}
Table~\ref{tab:cost_analysis} presents the cost comparison for evaluating the full corpus of 75 code files.

\begin{table}[h]
    \centering
    \caption{Cost comparison for evaluating 75 code files.}
    \label{tab:cost_analysis}
    \small
    \begin{tabular}{lrc}
        \toprule
        \textbf{Evaluator} & \textbf{Cost} & \textbf{Ratio} \\
        \midrule
        Claude 3.5 Sonnet & \$8.52 & 31.6$\times$ \\
        Claude 3.5 Haiku & \$1.66 & 6.1$\times$ \\
        \rowcolor{techBlue!10}
        \textbf{GPT-4o-mini} & \textbf{\$0.27} & \textbf{1.0$\times$} \\
        Human Experts (1--2 days) & \$200--400+ & 740--1480$\times$ \\
        \bottomrule
    \end{tabular}
\end{table}

Key findings from the cost analysis:
\begin{itemize}
    \item While Claude 3.5 Sonnet is the most accurate LLM, it costs \$8.52 per evaluation cycle, making it 31.6 times more expensive than GPT-4o-mini.
    \item \textbf{GPT-4o-mini} provides the best cost-to-accuracy ratio at only \$0.27 per cycle with correlations averaging $\rho \approx 0.75$, making it highly suitable for high-volume production pipelines.
    \item Human experts remain orders of magnitude more expensive (740--1480$\times$) while requiring 1--2 days per evaluation cycle, underscoring the economic case for LLM-assisted code review.
\end{itemize}

\subsection{Conclusion and Contributions}
The paper provides a systematic experimental comparison of LLM code review capabilities against human expert baselines.

Key contributions of the paper:
\begin{itemize}
    \item Establishes a rigorous parallel evaluation framework with statistically validated methodology (Shapiro-Wilk normality test, Spearman rank correlation).
    \item Demonstrates that Claude 3.5 Sonnet achieves the highest alignment with human expert consensus in 9 out of 12 code quality metrics, particularly in complex design principles.
    \item Reveals that model size strongly correlates with the ability to evaluate abstract software engineering concepts (SOLID principles).
    \item Provides a practical cost-benefit analysis showing that GPT-4o-mini offers the optimal balance between accuracy and cost for production-scale code review.
    \item Highlights that while LLMs cannot fully replace human reviewers, they can serve as effective first-pass screening tools, significantly reducing the workload on senior developers.
\end{itemize}
